\section{Unpenalized NegDRO}
\label{sec:unpenalized}

The unpenalized population objective
$\Phi(b)=\max_{w\in\Delta_m}\mathcal L(b,w)$ may be nonsmooth because the
inner objective is linear in $w$.  The stochastic extension does not attempt
to select a maximizing $w$ at every round.  It applies
\eqref{eq:common-primal-update}--\eqref{eq:common-dual-update} with $\mu=0$
and controls static regret against the same fixed witness $w^0$.

\subsection{Direction and comparator reduction}

Let $\varepsilon_{\rm id}$ be as in
\eqref{eq:penalized-identification-error}.
\begin{lemma}[Unpenalized fixed-witness direction]
\label{lem:unpenalized-fixed-witness}
For every $b\in\R^p$,
\begin{equation}
  (b-\beta^*)^\top\nabla_b\mathcal L(b,w^0)
  \ge
  \lambda\|b-\beta^*\|_2^2-\varepsilon_{\rm id}.
  \label{eq:unpenalized-fixed-witness-rsi}
\end{equation}
\end{lemma}
The proof is the same fixed-witness curvature and vector Young inequality used
in Lemma~\ref{lem:penalized-fixed-witness}.

Define
\begin{equation}
  \operatorname{Reg}_t^0
  :=\mathcal L(b_t,w^0)-\mathcal L(b_t,w_t).
  \label{eq:unpenalized-comparator-regret}
\end{equation}
The sign convention is comparator minus current, although the instantaneous
quantity may be negative.

\begin{lemma}[Unpenalized comparator reduction]
\label{lem:unpenalized-comparator-reduction}
For every round $t$,
\begin{equation}
  (b_t-\beta^*)^\top\nabla_b\mathcal L(b_t,w_t)
  \ge
  \lambda\|b_t-\beta^*\|_2^2
  -\varepsilon_{\rm id}-2\operatorname{Reg}_t^0.
  \label{eq:unpenalized-comparator-direction}
\end{equation}
\end{lemma}

\begin{proof}
The common term in the loss--gradient identity is independent of $w$, so
subtracting the identity at $w^0$ and $w_t$ gives
\[
  (b_t-\beta^*)^\top\nabla_b\mathcal L(b_t,w_t)
  =(b_t-\beta^*)^\top\nabla_b\mathcal L(b_t,w^0)
   -2\operatorname{Reg}_t^0.
\]
Lemma~\ref{lem:unpenalized-fixed-witness} completes the proof.
\end{proof}

The entropy mirror theorem described in Section~\ref{sec:penalized} no longer
takes Pinsker or the mirror-step conclusion as hypotheses.  Under its stated
simplex, positivity, and step-size conditions, it yields
\begin{equation}
  \sum_{t=1}^T\E[\operatorname{Reg}_t^0]
  \le
  \frac{\log m}{\eta_w}
  +\frac{\eta_wT G_{w,0}^2}{2}.
  \label{eq:unpenalized-static-regret}
\end{equation}
Substitution into the same expected projected recursion gives the final
unpenalized bounds.

\subsection{Verified unpenalized conclusions}

\begin{theorem}[Unpenalized finite-time expected squared error]
\label{thm:unpenalized-fixed-comparator}
Under the assumptions of Section~\ref{sec:setup}, with $\mu=0$, the recursion
satisfies
\begin{equation}
  \boxed{
  D_T
  \le
  \frac{\|b_{\rm Init}-\beta^*\|_2^2}
       {2\lambda\eta_bT}
  +\frac{\|v_{\rm SEM}\|_2^2}
       {\lambda^2(1+\gamma m)^2}
  +\frac{2\log m}{\lambda\eta_wT}
  +\frac{\eta_wG_{w,0}^2}{\lambda}
  +\frac{\eta_bG_b^2}{2\lambda}.}
  \label{eq:unpenalized-average-bound}
\end{equation}
\end{theorem}

\begin{theorem}[Unpenalized averaged-iterate expected squared error]
\label{thm:unpenalized-averaged-iterate}
The averaged iterate satisfies
\begin{equation}
  \boxed{
  \E\|\bar b_T-\beta^*\|_2^2
  \le
  \frac{\|b_{\rm Init}-\beta^*\|_2^2}
       {2\lambda\eta_bT}
  +\frac{\|v_{\rm SEM}\|_2^2}
       {\lambda^2(1+\gamma m)^2}
  +\frac{2\log m}{\lambda\eta_wT}
  +\frac{\eta_wG_{w,0}^2}{\lambda}
  +\frac{\eta_bG_b^2}{2\lambda}.}
  \label{eq:unpenalized-averaged-iterate-bound}
\end{equation}
\end{theorem}

The identification contribution in the direction inequality is
$\varepsilon_{\rm id}$; division by $\lambda$ in the primal recursion produces
the displayed final bias
$\|v_{\rm SEM}\|_2^2/[\lambda^2(1+\gamma m)^2]$.

Define the exact inverse-square-root coefficient
\begin{equation}
  C_0
  :=\frac{\|b_{\rm Init}-\beta^*\|_2^2}{2\lambda}
    +\frac{2\log m}{\lambda}
    +\frac{G_{w,0}^2}{\lambda}
    +\frac{G_b^2}{2\lambda}.
  \label{eq:unpenalized-rate-constant}
\end{equation}

\begin{corollary}[Unpenalized inverse-square-root rate]
\label{cor:unpenalized-inv-sqrt}
If $\eta_b=\eta_w=1/\sqrt T$, then
\begin{equation}
  D_T\le B_{\rm id}+\frac{C_0}{\sqrt T}.
  \label{eq:unpenalized-inv-sqrt-rate}
\end{equation}
\end{corollary}

This public rate theorem controls $D_T$, not the terminal state $b_{T+1}$.
The separate
averaged-iterate theorem retains the general positive step sizes and the same
finite-time right-hand side; no additional averaged inverse-square-root public
entry point is asserted here.
